\section*{Capítulo \ref{ch:g10:storms-and-cyclones} --- Trovoadas e ciclones tropicais}

\begin{solution}{exo:g10:storms-and-cyclones:1}
Uma \omterm{def:g10:storms-and-cyclones:thunderstorm}{trovoada} é uma tempestade produzida por um cumulonimbo, com \omterm{def:g10:storms-and-cyclones:lightning}{raios} e
trovões, em geral aguaceiros fortes e rajadas, às vezes granizo. O \omterm{def:g10:storms-and-cyclones:lightning}{raio} é
uma faísca elétrica gigante que descarrega as regiões carregadas da
tempestade. O trovão é o som do ar do canal do \omterm{def:g10:storms-and-cyclones:lightning}{raio}, aquecido tão de repente
que se expande num choque.
\end{solution}

\begin{solution}{exo:g10:storms-and-cyclones:2}
Um \omterm{def:g10:storms-and-cyclones:tropical-cyclone}{ciclone tropical} é uma tempestade com centenas de quilômetros de
diâmetro, formada sobre um mar tropical quente em torno de um centro de
baixa pressão, com ventos que espiralam para dentro e faixas de \omterm{def:g10:storms-and-cyclones:thunderstorm}{trovoadas}.
O \omterm{def:g10:storms-and-cyclones:tropical-cyclone}{olho} é o seu centro calmo e quase sem nuvens; a \omterm{def:g10:storms-and-cyclones:tropical-cyclone}{parede do olho} é o anel de
\omterm{def:g10:storms-and-cyclones:thunderstorm}{trovoadas} altíssimas em volta do \omterm{def:g10:storms-and-cyclones:tropical-cyclone}{olho}. Os ventos são mais fortes na \omterm{def:g10:storms-and-cyclones:tropical-cyclone}{parede
do olho}.
\end{solution}

\begin{solution}{exo:g10:storms-and-cyclones:3}
\num{70}~nós: categoria 1. \num{120}~nós: categoria 4.
\num{140}~nós: categoria 5. Os dois últimos são furacões intensos
(categoria 3 ou mais).
\end{solution}

\begin{solution}{exo:g10:storms-and-cyclones:4}
$343 \times 9 \approx \qty{3.1}{km}$, depois $343 \times 15 \approx
\qty{5.1}{km}$. O segundo clarão está mais longe: a tempestade está se
afastando.
\end{solution}

\begin{solution}{exo:g10:storms-and-cyclones:5}
O \omterm{def:g10:storms-and-cyclones:lightning}{raio} tende a cair no ponto mais alto das redondezas: uma árvore sozinha
num campo é esse ponto, e a corrente de uma descarga desce por ela e se
espalha pelo chão em volta. Entre num prédio ou num carro fechado e fique ali
até bem depois do último trovão.
\end{solution}

\begin{solution}{exo:g10:storms-and-cyclones:6}
Um nó vale $1852/3600 \approx \qty{0.514}{m/s}$. \num{64}~nós: $64 \times
0.514 \approx \qty{32.9}{m/s}$, e $32.9 \times 3.6 \approx
\qty{119}{km/h}$. \num{137}~nós: $\qty{70.5}{m/s}$, cerca de
\qty{254}{km/h} (a tabela oficial, arredondada passando por outra unidade,
dá \qty{252}{km/h}).
\end{solution}

\begin{solution}{exo:g10:storms-and-cyclones:7}
$44 \times 86\,400 \approx \num{3.8e6}$ \omterm{def:g10:storms-and-cyclones:lightning}{raios} por dia, e $3.8\times10^6
\times 365.25 \approx \num{1.4e9}$, cerca de 1.4 bilhão por ano.
\end{solution}

\begin{solution}{exo:g10:storms-and-cyclones:8}
$\Delta p = 1010 - 935 = \qty{75}{hPa}$: uma elevação de cerca de
\qty{75}{cm}.
\end{solution}

\begin{solution}{exo:g10:storms-and-cyclones:9}
Eles se formam sobre os oceanos tropicais quentes, sobretudo entre
\qty{10}{\degree} e \qty{30}{\degree} de latitude nos dois hemisférios, com
muito poucos no Atlântico Sul. Derivam para oeste e depois se curvam em
direção aos polos, onde se enfraquecem sobre águas mais frias. Estão
ausentes do equador porque ali a rotação da Terra não dá nenhum giro ao ar
que entra: sem ele, o ar flui direto para a baixa pressão e a preenche, e
nenhuma tempestade em rotação pode começar.
\end{solution}

\begin{solution}{exo:g10:storms-and-cyclones:10}
$1031 - 1.56 \times 70 \approx \qty{922}{hPa}$; $1031 - 1.56 \times 40
\approx \qty{969}{hPa}$.
\end{solution}

\begin{solution}{exo:g10:storms-and-cyclones:11}
Sobre a terra, a tempestade perde o seu suprimento de água que evapora, e
com ele o \omterm{def:g10:earth-energy-budget:latent-heat}{calor latente} que a move; a superfície rugosa também freia os
ventos. Sobre águas frias, a evaporação é muito mais fraca, já que o ar frio
contém muito menos vapor de água na saturação: o suprimento de \omterm{def:g10:earth-energy-budget:latent-heat}{calor latente}
diminui.
\end{solution}

\begin{solution}{exo:g10:storms-and-cyclones:12}
A \qty{0}{\degreeCelsius}: $\sqrt{1.4 \times 8.314 \times 273.15/0.02897}
\approx \qty{331}{m/s}$; a \qty{30}{\degreeCelsius}, com $T =
\qty{303.15}{K}$: \qty{349}{m/s}. Para \qty{9}{s}: $9 \times 331 \approx
\qty{2980}{m}$ e $9 \times 349 \approx \qty{3140}{m}$, uma diferença de
cerca de \qty{160}{m}, cerca de \qty{5}{\percent}. A regra dá \qty{3}{km}:
o bastante para saber a que distância está a tempestade, que é tudo o que se
pede dela.
\end{solution}

\begin{solution}{exo:g10:storms-and-cyclones:13}
$\Delta p = 1010 - 940 = \qty{70}{hPa}$: a elevação barométrica é de cerca
de \qty{0.7}{m}, a parte do vento $4.0 - 0.7 = \qty{3.3}{m}$, mais de quatro
vezes mais. Na maré alta, a água chega a $4.0 + 1.5 = \qty{5.5}{m}$ acima do
nível médio do mar.
\end{solution}

\begin{solution}{exo:g10:storms-and-cyclones:14}
\omterm{def:g10:storms-and-cyclones:tornado}{Tornado}: uma volta mede $2\pi \times 100 \approx \qty{628}{m}$, percorrida
em $628/90 \approx \qty{7}{s}$. Furacão: $2\pi \times 30\,000 \approx
\qty{188}{km}$, percorridos em $188\,000/60 \approx \qty{3140}{s}$, cerca de
\qty{52}{min}. O \omterm{def:g10:storms-and-cyclones:tornado}{tornado} gira umas 450 vezes mais depressa: é um redemoinho
pequeno e muito rápido; o furacão, um redemoinho imenso e mais lento.
\end{solution}

\begin{solution}{exo:g10:storms-and-cyclones:15}
$24/3888 \approx \qty{0.6}{\percent}$; $414/3888 \approx
\qty{10.6}{\percent}$; $121/3888 \approx \qty{3.1}{\percent}$. Se a latitude
não importasse, cerca de \qty{9}{\percent} das tempestades, umas 350,
começariam a menos de \qty{5}{\degree} do equador; só 24 o fazem, cerca de
quinze vezes menos. Perto do equador, a rotação da Terra é fraca demais para
dar início ao giro.
\end{solution}

\begin{solution}{pb:g10:storms-and-cyclones:1}
\textbf{1.} $\pi \times (665\,000)^2 \approx \qty{1.39e12}{m^2}$.

\textbf{2.} $0.015 \times 1.39\times10^{12} \approx \qty{2.08e10}{m^3}$.

\textbf{3.} Um quilômetro cúbico são \qty{e9}{m^3}: cerca de \qty{21}{km^3}.

\textbf{4.} $2.08\times10^{10} \times 1000 \approx \qty{2.08e13}{kg}$.

\textbf{5.} $2.08\times10^{13} \times 2.5\times10^6 \approx \qty{5.2e19}{J}$.

\textbf{6.} $5.2\times10^{19}/86\,400 \approx \qty{6.0e14}{W}$.

\textbf{7.} $6.0\times10^{14}/1.22\times10^{17} \approx 0.005$, cerca de
\qty{0.5}{\percent}.

\textbf{8.} $240 \times 1.39\times10^{12} \approx \qty{3.3e14}{W}$. A
tempestade libera quase o dobro do calor da luz solar absorvida por toda a
sua área.

\textbf{9.} A tempestade não funciona com a luz solar que cai sobre ela: ela
libera um calor armazenado como vapor de água, evaporado de uma área maior
de mar quente e trazido pelos ventos. A tempestade concentra uma energia
reunida sobre uma região muito maior.

\textbf{10.} A mesma massa que a chuva, cerca de \qty{2.1e13}{kg} por dia (o
vapor trazido precisa ser renovado tão depressa quanto cai como chuva).

\textbf{11.} $6.0\times10^{14}/82 \approx \qty{7.3e12}{m^2}$, cerca de cinco
vezes a área da tempestade: a tempestade precisa tirar vapor de uma área
muito maior, e os seus ventos fazem o mar sob ela evaporar muito mais
depressa que a média.

\textbf{12.} \qty{24.1}{g/kg} a \qty{28}{\degreeCelsius} contra
\qty{14.7}{g/kg} a \qty{20}{\degreeCelsius}: $24.1/14.7 \approx 1.64$, cerca
de \qty{64}{\percent} a mais.

\textbf{13.} A água quente evapora depressa e o ar quente contém muito mais
vapor: o ar que entra traz muito mais \omterm{def:g10:earth-energy-budget:latent-heat}{calor latente} para liberar e,
saturado, sobe livremente. Sobre mares mais frios, o suprimento é pequeno
demais para mover o motor.

\textbf{14.} Sobre a terra não há mar para evaporar: o suprimento de \omterm{def:g10:earth-energy-budget:latent-heat}{calor
latente} para, e o chão rugoso freia os ventos.

\textbf{15.} $1.5\times10^{12} \times 86\,400 \approx \qty{1.3e17}{J}$.

\textbf{16.} $96 \times 1852/3600 \approx \qty{49.4}{m/s}$.

\textbf{17.} $\frac12 \times 1.16 \times 49.4^2 \approx \qty{1.4e3}{J}$.

\textbf{18.} $5.2\times10^{19}/1.3\times10^{17} = 400$: o \omterm{def:g10:earth-energy-budget:latent-heat}{calor latente} é
cerca de 400 vezes a energia fornecida aos ventos.

\textbf{19.} A maior parte dele aquece o ar alto da tempestade, que se
espalha pelo topo e cede o calor ao espaço como radiação infravermelha; os
próprios ventos perdem a sua energia por atrito contra o mar e o chão.

\textbf{20.} $1/400 = \qty{0.25}{\percent}$ (o programa de cálculo dá uma
razão de 402): \textbf{cerca de uma parte em 400 do calor liberado pela
condensação se torna vento}.
\end{solution}
